% !TeX root = ./../main.tex

\chapter{Przegląd literatury}\label{chapter:literature_review}

\agentinfo{
Cel rozdziału: ustalić, w jakim stopniu istniejące metody korekcji, mitygacji i redukcji błędów pozwalają dziś prowadzić kwantowe przetwarzanie i klasyfikację obrazów na urządzeniach ery NISQ, oraz wyprowadzić lukę badawczą, którą wypełnia praca.
Narracja prowadzi od metod ogólnych (korekcja, mitygacja), przez metody oparte na pomiarze i uczeniu maszynowym, charakteryzację szumu, do zastosowań obrazowych i klasyfikacji, a kończy się syntezą i luką.
Definicje metod znajdują się w zapleczu teoretycznym \agentref{chapters/theoretical_background_03_QEC.tex}; tutaj analizowane są prace: zakres, warunki eksperymentów, sprzęt lub symulator, miary, wyniki, koszty, ograniczenia.

Docelowa objętość rozdziału: około 14 stron.
Punkt wyjścia literaturowy: prace cytowane w artykułach autora (bibliography/biblio.bib).
Literatura z artykułów kończy się w dużej mierze na latach 2021--2023; przegląd wymaga uzupełnienia o nowsze prace, w szczególności o eksperymentalne demonstracje korekcji błędów poniżej progu, nowsze przeglądy mitygacji błędów oraz prace o uczeniu maszynowym w mitygacji.
Każdą nową pozycję zweryfikować w źródle (DOI, tytuł, autorzy) przed dodaniem do bazy; nie tworzyć wpisów z pamięci.

Zasady analizy:
rozróżniać deklaracje autorów, wyniki poparte eksperymentem lub dowodem oraz własne wnioski;
dla istotnych twierdzeń notować lokalizację w źródle (sekcja, tabela, wzór);
brak informacji w publikacji oznaczać jako ,,nie podano'';
brak wyników wyszukiwania nie dowodzi istnienia luki badawczej.
}

\section{Cel i zakres przeglądu}\label{sec:lr:scope}
\agentinfo{
Wstęp do rozdziału: zakres tematyczny i czasowy, kryteria doboru prac, przebieg argumentacji.
Uzupełnić na końcu, po opracowaniu pozostałych części.
}

\section{Kwantowa korekcja błędów}\label{sec:lr:qec}
\agentinfo{
Stan badań nad kodami stabilizatorowymi i powierzchniowymi oraz ich realizacjami eksperymentalnymi.
Prace cytowane we wstępie artykułu ICCS \agentref{artykuły/Quantum_Errors_ICCS/main.tex}: aaronson_improved_2004, nguyen_quantum_2018, nguyen_novel_2020, lv_explicit_2020, ryan-anderson_realization_2021, dymarsky_quantum_2021, bravyi_correcting_2018, litinski_game_2019.
Uzupełnić o nowsze demonstracje sprzętowe i o analizę narzutu zasobów.
Wniosek, do którego prowadzi sekcja: pełna korekcja pozostaje kosztowna dla obliczeń wymagających całej dostępnej objętości kwantowej, takich jak przetwarzanie obrazów.
}

\section{Mitygacja błędów w erze NISQ}\label{sec:lr:qem}
\agentinfo{
Ekstrapolacja do zerowego szumu, probabilistyczne odwracanie szumu, mitygacja błędów pomiaru, weryfikacja symetrii, metody oparte na regresji danych z obwodów Clifforda.
Prace wyjściowe: endo_practical_2018, maciejewski_mitigation_2020, richardson-extrapolation, quasi-probability.
Dla każdej grupy metod: założenia o szumie, koszt próbkowania, skala eksperymentów, ograniczenia.
Szczególnie dokładnie omówić mitygację błędów pomiaru i metody kalibracyjne, bo są najbliższe PDU.
}

\section{Redukcja błędów oparta na pomiarze i uczeniu maszynowym}\label{sec:lr:ml}
\agentinfo{
Metody korygujące wynik po stronie klasycznej: modele uczone na parach wynik zaszumiony--wynik idealny, sieci neuronowe do mitygacji, odszumianie rozkładów wyników, modele generatywne.
Także dekodery neuronowe dla kodów korekcyjnych, jako kontekst.
Sekcja ma pokazać, czy i jak wcześniej stosowano modele obraz--obraz do wyników obliczeń kwantowych; jest to tło dla rozdziału 6.
}

\section{Charakteryzacja i modelowanie szumu urządzeń}\label{sec:lr:noise}
\agentinfo{
Tomografia procesu, gate set tomography, randomized benchmarking, uczenie modeli szumu, szum niemarkowowski, dryf w czasie, przesłuch.
Stabilność błędu między kalibracjami: prace potwierdzające lub podważające założenie metody PDU.
Symetrie a szum i korekcja błędów.
Źródło: sekcja Related Work artykułu \agentref{artykuły/quantum_noise_elsevier_submission/main.tex} oraz jego bibliografia.
Jest to tło dla rozdziałów 7 i 9.
}

\section{Kwantowe przetwarzanie obrazów w obecności szumu}\label{sec:lr:qimg}
\agentinfo{
Przeglądy metod reprezentacji obrazów (na przykład cai_survey_2018), realizacje na rzeczywistych urządzeniach, badania odporności reprezentacji na szum, maksymalne rozmiary obrazów osiągane eksperymentalnie.
Definicje FRQI, NEQR i LPIQE są w zapleczu teoretycznym \agentref{chapters/theoretical_background_05_QIMG.tex}; tutaj porównanie wyników i warunków eksperymentów.
Uwzględnić także uczone i przybliżone kodowania obrazów (autoenkodery kwantowe i hybrydowe, kodowania wariacyjne), które są tłem dla rozdziału 7.
Jest to tło dla rozdziałów 5 i 7.
}

\section{Klasyfikacja obrazów z użyciem komputerów kwantowych a szum}\label{sec:lr:classification}
\agentinfo{
Klasyfikatory kwantowe i hybrydowe dla obrazów, wpływ szumu na trening i wnioskowanie, płaskowyże gradientu indukowane szumem, trening uwzględniający szum, szum jako regularyzator lub zasób.
Sekcja łączy tło dla rozdziału 8 (wpływ szumu i redukcji błędów na klasyfikację) i rozdziału 9 (szum jako dodatkowy stopień swobody w uczeniu).
Uwzględnić pracę współautorską o klasyfikacji obrazów FRQI kwantową siecią rekurencyjną \agentref{bib:potempa_frqi_pairs_2025}.
Źródło części o strukturze grupowej: sekcja Related work artykułu \agentref{artykuły/noisy_lie_group_learners/main.tex}.
}

\section{Synteza stanu badań i luka badawcza}\label{sec:lr:gap}
\agentinfo{
Lukę wyprowadzić z porównania prac, a nie z deklaracji.
Kandydaci na elementy luki, do sprawdzenia w literaturze wraz z możliwymi kontrprzykładami:
(a) brak metod redukcji błędów zorientowanych na problem, które nie zużywają kubitów, bramek ani dodatkowych wykonań obwodu poza kalibracją;
(b) brak systematycznego porównania utraty danych w reprezentacjach obrazów na rzeczywistych urządzeniach;
(c) brak oceny modeli generatywnych obraz--obraz jako etapu redukcji błędu po korekcie kalibracyjnej;
(d) brak oceny wpływu redukcji błędów na klasyfikację obrazów kodowanych kwantowo;
(e) ograniczona wiedza o strukturze szumu wykraczającej poza modele zachowujące symetrię.
Każdy element luki opatrzyć uzasadnieniem źródłowym, zakresem niepewności i odwołaniem do rozdziału, który go adresuje.
Warto zakończyć tabelą: metoda, dodatkowe kubity, dodatkowe bramki, dodatkowe wykonania, założenia, zastosowanie do obrazów.
}
